% id: 66-35
% book: 베이직쎈 확률과 통계 · 04 조건부확률
% section: 개념 19 $\mathrm{P}(B)=\mathrm{P}(A\cap B)+\mathrm{P}(\comp{A}\cap B)$를 이용한 확률
% passage: 다음에 답하시오.
% figure: none
% answer: 
% answer_source: 
% variant_level: 0 (원본 전사)
% 이 파일은 build.mjs 가 items.json 에서 만든다. 고칠 때는 items.json 을 고친다.
\smash{\raisebox{1.5mm}{\dmkichul{베이직쎈 확률과 통계 66쪽 35번}}}
노란 공 4개, 검은 공 6개가 들어 있는 $\mathrm{A}$ 주머니와 노란 공 2개, 검은 공 4개가 들어 있는 $\mathrm{B}$ 주머니 중 임의로 하나를 택하여 한 개의 공을 꺼낼 때, $\mathrm{A}$ 주머니를 택하는 사건을 $A$, $\mathrm{B}$ 주머니를 택하는 사건을 $B$, 꺼낸 공이 검은 공인 사건을 $E$라 하자. 다음을 구하시오.
\dmsubprob{1} 꺼낸 공이 검은 공일 확률\dmhint{$\mathrm{P}(A)=\dfrac{1}{2}$, $\mathrm{P}(B)=\blank$이므로\\ $\mathrm{P}(E)=\mathrm{P}(A\cap E)+\mathrm{P}(\blank[4em])$\\ $\hphantom{\mathrm{P}(E)}=\mathrm{P}(A)\mathrm{P}(E\mid A)+\mathrm{P}(B)\mathrm{P}(\blank[3.2em])$\\ $\hphantom{\mathrm{P}(E)}=\dfrac{1}{2}\cdot\dfrac{\blank}{10}+\dfrac{1}{2}\cdot\dfrac{\blank}{6}=\blank$}
\dmsubprob{2} 꺼낸 공이 검은 공일 때, 택한 주머니가 $\mathrm{A}$일 확률\dmhint{$\mathrm{P}(A\mid E)=\dfrac{\mathrm{P}(A\cap E)}{\mathrm{P}(\blank)}=\dfrac{\blank}{\blank}=\blank$}
